\section{Relation to the Note}


本節では、本稿と note 本編との関係を整理する。

本稿は、note 本編を置き換えるためのものではない。

note 本編は、意識についての直観的な読み物であり、体験、比喩、違和感、日常的な感覚から出発する層である。

それに対して本稿は、その背後にある構造を固定するための理論基盤である。

したがって、本稿の役割は、note 本編を理論で上書きすることではなく、note 本編で直観的に示されていたものを、時間的非閉包、速度スケール差、ログ参照、意味勾配という構造の中へ再配置することである。

\subsection{The Note as Intuitive Layer}

note 本編では、読者が最初に出会うべきなのは数式ではない。

読者が最初に出会うべきなのは、自分の意識経験に対する違和感である。

なぜ、自分が決めたと思う前に身体は動き始めているのか。

なぜ、言葉にした瞬間に感覚は少し変質するのか。

なぜ、強い記憶や感情は何度も意識に戻ってくるのか。

なぜ、意識は自分の中にあるように感じられるのに、探すと場所が見つからないのか。

note 本編は、これらの問いを直観的に扱う。

そこでは、理論用語は最小限でよい。

重要なのは、読者が「意識は何かが生成されているというより、何かが開いているのかもしれない」と感じられることである。

\subsection{The Present Paper as Structural Layer}

一方、本稿は、note 本編で提示された直観を構造化する。

本稿では、意識を以下のように定式化した。

\begin{itemize}
\item 意識は生成物ではなく、ログ参照状態である。
\item 意識は単一部位ではなく、時間的非閉包レジームである。
\item 意識は高速層と低速層の速度差 $\Delta v$ に依存する。
\item 意識には結合粘性 $\mu$ と内部遅延 $\tau$ が必要である。
\item ログ参照は意味勾配 $\nabla \Phi$ によって方向づけられる。
\item 有効トルク様量 $\mathcal{T}_{\Phi}$ は、速度差と意味勾配の結合を表す。
\item クオリアは意識の本体ではなく、沈殿境界のライブな前景である。
\item 測定はそのライブな境界を残余へ変換する。
\item 意識は局在せず、複数の局所基盤が特定の時間的関係に入ったときに開かれる。
\end{itemize}

これらは、note 本編で前面に出す必要はない。

しかし、note 本編の背後にある構造として保持される必要がある。

\subsection{What Should Flow Back to the Note}

本稿から note 本編へ還流させるべきものは、すべての数式ではない。

還流させるべきなのは、読者の理解を助ける最小限の構造である。

具体的には、次の要素である。

\begin{enumerate}
\item 意識は「作られる」のではなく「開かれる」という表現。
\item 高速な処理と低速な記憶沈殿のあいだに意識が開くという直観。
\item 自我は主体ではなく、速度差の変換点であるという整理。
\item クオリアは保存されるものではなく、沈殿境界のライブな前景であるという整理。
\item 測定や内省は、クオリアをそのまま捕まえるのではなく、ログ化するという整理。
\item 意識は脳内の一点ではなく、時間的条件として成立するという整理。
\end{enumerate}

これらは、note 本編の読者にとって十分に直観的でありながら、本稿の理論基盤とも接続している。

\subsection{What Should Not Flow Back Directly}

逆に、本稿から note 本編へ直接持ち込むべきでないものもある。

第一に、数式を過剰に持ち込むべきではない。

$\Delta v$、$\mu$、$\tau$、$\mathcal{T}_{\Phi}$、$\nabla \Phi$ は、本稿では重要な構造変数である。

しかし、note 本編では、それらをすべて明示すると、読み物としての流れが崩れる。

第二に、病理的逸脱を診断のように扱うべきではない。

本稿の pathologies は、疾患分類ではなく、consciousness window からの逸脱方向である。

note 本編でこれを扱う場合も、医学的説明ではなく、意識状態の変形として慎重に扱う必要がある。

第三に、IFGT/VED の用語を前面に出しすぎるべきではない。

note 本編の役割は、理論名を理解させることではなく、読者に構造を体験させることである。

理論は前に出るのではなく、裏から支えるべきである。

\subsection{Layered Translation}

本稿と note 本編の関係は、単純な要約ではない。

それは、層をまたいだ翻訳である。

L3 では、意識は次のように書かれる。

\begin{equation}
\text{Consciousness}
\sim
\text{Log Referencing}
\left(
\Delta v,
\mu,
\tau,
\mathcal{T}_{\Phi},
\nabla \Phi
\right)
\end{equation}

しかし、note 本編では、これをそのまま出す必要はない。

note 本編では、次のように言い換えればよい。

意識とは、速く流れる処理と、ゆっくり沈む記憶のあいだに、一時的に開く通路である。

この翻訳が重要である。

同じ構造を、異なる解像度で示す。

これが、note 本編と本稿の関係である。

\subsection{Relation to Appendices}

note 本編と本稿のあいだには、補遺や中間説明層が置かれる。

本稿は L3 として、構造を固定する。

補遺は L2 として、その構造を読者に説明する。

必要であれば、L2.5 として、数式の意味を直観へ戻す橋渡しを置く。

note 本編は L1 として、読者が最初に触れる体験層である。

この層構造を混同しないことが重要である。

L3 が L1 を支配してはならない。

L1 が L3 を曖昧にしてもならない。

それぞれの層は異なる役割を持つ。

本稿の目的は、L1 の文章を硬くすることではなく、L1 が安心して柔らかく書けるように、背後の構造を固定することである。

\subsection{Summary}

本稿と note 本編の関係は、置換ではなく再配置である。

\begin{itemize}
\item note 本編は、体験と直観の層である。
\item 本稿は、構造と方程式の層である。
\item 本稿の内容は、すべてを note 本編へ持ち込む必要はない。
\item note 本編へ還流すべきなのは、読者の直観を助ける最小限の構造である。
\item 理論用語や数式は、必要な場合に補遺や中間層で支える。
\item L3 は L1 を上書きしない。
\item L3 は L1 を裏から支える。
\end{itemize}

したがって、本稿は note 本編の完成版ではない。

本稿は、note 本編が読み物として自由に呼吸するための、背後の骨格である。
